\section*{序論}
10月28日朝までの生成AI関連Xでは、操作可能な世界モデルと小規模オープンウェイトモデルの進展が目立った。Google DeepMindのGenie 4はリアルタイム操作と音声を伴うシミュレーションを掲げ、Tencent HunyuanのHy-Nomasは単一GPUで動かせる公開モデルとして対照的に紹介された。Liquid AIのLFM3も、端末側でのマルチモーダル処理を志向する動きとして注目された。

研究面では、OpenAIとAnthropicが指数時間仮説（ETH）の反証を得たとする主張が議論を呼んだ。ただし、数学的な正当性は独立検証が終わっておらず、確定した成果として扱うことはできない。Cloudflareの過去のデータ露出インシデントの開示も、AI推論サービスの隔離設計を考える材料となった。

本号ではGrokの8件を収集順に整理する。公式発表が観測窓外にある場合や、ベンダー・コミュニティの主張にとどまる場合は、その境界を明記する。独立した一次資料の再検証は行っていない。
